\def\chaptercode{LP11}    % Code du chapitre
\def\chaptername{Présentation de l'outil Capytale}        % Nom du chapitre

\titre{}

\begin{methode}{Accéder à Capytale depuis l'ent}{}
	\smallskip
	\begin{minipage}{0.59\linewidth}
		\begin{enumerate}
			\item s'identifier sur l'ent via educonnect.
			\item dans le menu de gauche: \verb+ Pédagogie / Capytale+{}
			\item entrer le code dans \verb+Accéder à une activité+{} puis \verb+GO !+{}
			%\item cliquer sur \verb+Go!+
		\end{enumerate}
	\end{minipage}\hfill
	\begin{minipage}{0.39\linewidth}
		\includegraphics[width=1\linewidth]{capytale1.jpg}
	\end{minipage}
	\smallskip
\end{methode}

\begin{methode}{Accéder à Capytale via un lien direct}{}
	\smallskip
	\begin{enumerate}
		\item cliquer sur le lien fourni par l'enseignant \verb+ https://capytale2.ac-paris.fr/web/c/xxxx-xxxxxxx + 
		\item s'identifier avec les identifiants de l'ent
		%\item vous arrivez sur la page de Capytale.
	\end{enumerate}
\end{methode}

\begin{methode}[breakable, enhanced jigsaw]{Prise en main de la fenêtre Capytale}{}
	\center\includegraphics[width=0.7\linewidth]{zonesCapytale.jpg}


	\begin{liste}
		\item pensez à sauvegarder régulièrement. La sauvegarde automatique est parfois aléatoire.
		\item les cellules peuvent être de 4 types: \texttt{Code, Markdown, Texte brut, Titre}.
		\begin{itemize}
			\item le type \texttt{Texte Brut} est uniquement du texte sans aucune mise en forme.
			\item le type \texttt{Markdown} est du texte avec une mise en page légère. %Vous pouvez cliquer sur une des cellules au type \texttt{Markdown} pour voir ce qui se cache derrière.
			\item le type \texttt{Code} est une cellule contenant en général du code \texttt{Python} que l'on peut exécuter
			
		\end{itemize}

		%\newpage
		\item vous pouvez créer vous même des cellules si besoin en choisissant son type.
		\item l'ordre d'exécution est important. Avant toute exécution de code, il n'y a aucun nombre entre les crochets: \texttt{[ ]}. Une fois exécutée, le nombre entre crochet va s'incrémenter.
		%\item la gestion des fichiers attachés permet de pouvoir insérer et créer des fichiers comme des photos, du texte etc.
		
		\end{liste}
	\end{methode}


\begin{methode}{Fonctionnement de l'outil Capytale}{}
	\smallskip
	\begin{minipage}{0.2\linewidth}
		\includegraphics[width=1\linewidth]{acceder.jpg}
	\end{minipage}
	\begin{minipage}{0.7\linewidth}
		Lorsque vous ouvrez une activité sous Capytale, vous en faites techniquement une copie qui vous appartient. Vous pourrez y accéder de nouveau les séances suivantes en cliquant sur \texttt{Accéder à vos activités}.
	\end{minipage}
	
	\begin{minipage}{0.2\linewidth}
		\includegraphics[width=1\linewidth]{rendreTravail.jpg}
	\end{minipage}
	\begin{minipage}{0.7\linewidth}
		Si vous utilisez Capytale pour faire un devoir alors vous devez, à la fin de votre séance cliquer sur \texttt{Rendre ce travail}.
	\end{minipage}
\end{methode}

%\begin{tcolorbox}
%	\begin{liste}
%		\item Tutoriel accès Capytale élève:
%		
%		\href{https://tube-sciences-technologies.apps.education.fr/w/uPcXUyaEMuSyHf3h5WsL5M}{https://tube-sciences-technologies.apps.education.fr/w/uPcXUyaEMuSyHf3h5WsL5M}
%		
%		\item Tutoriel mise en page Markdown:
%		
%		\href{https://www.markdowntutorial.com/fr/}{https://www.markdowntutorial.com/fr/}
%	\end{liste}
%\end{tcolorbox}


\resetalltcb{}
\newpage





